\documentclass{article}

\usepackage{graphicx}
\usepackage{booktabs}

\sloppy

\title{Latex Curso}
\author{Geovana Lima, Reuben Claudino, Victor Pessoa Oliveira Ortins}

\date{Março de 2026}

\begin{document}

\maketitle

\section{Introdução ao Documento Básico}

Este é o nosso primeiro documento criado do zero. A estrutura básica de um arquivo \LaTeX{} é dividida em duas partes: o preâmbulo e o corpo do texto. Começamos sempre com a definição da classe do documento, que neste caso é um artigo. 

O cabeçalho e as informações iniciais, como título, autores e data, foram configurados no preâmbulo. Para que essas informações apareçam renderizadas na página, utilizamos um comando específico logo após iniciar o documento.

\section{Parágrafos e Formatação de Texto}

Escrever no \LaTeX{} exige entender como ele processa o texto. Para criar um novo parágrafo, não basta apertar "Enter" uma vez; é necessário deixar uma linha completamente em branco no código.

Também podemos aplicar estilos diretos ao texto para destacar informações. Por exemplo, podemos deixar uma palavra em \textbf{negrito} ou aplicar o estilo \textit{itálico}. O compilador se encarrega de aplicar a formatação sem que precisemos selecionar o texto com o mouse.

\section{Matemática e Equações}

O grande poder desta ferramenta está na escrita matemática. Podemos inserir fórmulas no meio do texto, colocando os elementos entre cifrões. Um exemplo clássico é a equação da equivalência massa-energia: $E = mc^2$.

Se precisarmos que a equação ganhe destaque e seja numerada, utilizamos ambientes específicos, como o ambiente de equação. Veja o exemplo de uma função polinomial abaixo:

\begin{equation}
f(x) = x^2 + 2x - 3
\end{equation}

Dessa forma, construímos estruturas complexas mantendo uma formatação padronizada e profissional.

\section{Imagens}

Para colocar imagens em um documento Latex utilizamos a tag \textit{figure}. Nela, podemos definir a escala da imagem (seu tamanho em comparação ao tamanho da página), escolher o arquivo da imagem e colocar a legenda

\begin{figure}[h]
  \centering
  \includegraphics[scale=0.1]{Flamengo-RJ.png}
  \caption{Uma vez Flamengo, sempre Flamengo.}
\end{figure}

\section{Tabelas}

A fim de inserir tabelas no latex, é preciso utilizar a tag \textit{tabular} e preencher os dados da tabela

\begin{table}[h]
    \centering
    \caption{Exemplo}
    \label{Tabela_Flamengo2}
    \begin{tabular}{llcc}
      \toprule
      Atleta & Gols & Assistências \\ \hline
      Arrascaeta & 12 & 5 \\
      ER7 & 8 & 10 \\ \hline
    \end{tabular}
\end{table}

\begin{table}[h]
    \centering
    \caption{Artilheiros do Flamengo}
    \label{Tabela_Flamengo}
    \begin{tabular}{llcc}
      \toprule
      Atleta & Gols & Assistências \\ \hline
      Gabigol & 12 & 5 \\
      Bruno Henrique & 8 & 10 \\ \hline
    \end{tabular}
\end{table}

\section{Referências}

Podemos também referenciar elementos anteriores do texto, se eles foram identificados pela tag \textit{label}. Caso tenham sido, é só utilizar a tag \textit{ref} e referenciar o label textual dado pelo escritor. Não há necessidade, então, de atualizar manualmente os números de tabelas/imagens ao inserir novos elementos durante o texto. Segue exemplo abaixo:

Como mostrado na tabela \ref{Tabela_Flamengo}, Gabigol é o artilheiro do Flamengo com 12 gols


\section{Bibliografia e Sumário}

É sempre possível adicionar citações ao seu documento. Idealmente, uma citação é feita no meio do texto e no final do artigo ela será devidamente adicionada. Para isso, crie um arquivo chamado referencias.bib e centralize lá suas referências. Uma citação é feita utilizando a tag \textit{cite}. É preciso também adicionar uma tag \textit{bibliography} com o nome do arquivo criado para centralizar as referências (nesse caso, referencias).

Referências de artigos podem ser facilmente exportadas na maioria dos sites agregadores de artigos. É só colocar no arquivod e referencias que ela já pode ser utilizada

O detector de faces RetinaFace foi introduzido inicialmente no artigo \cite{RetinaFacePaper:2019}, em que este modelo atingiu o estado de arte no que tange a área de detecção de faces

É possível também adicionar um sumário onde desejado no texto

\tableofcontents


\bibliographystyle{plain}
\bibliography{referencias}




\end{document}

